% DRAFT 5 — Expanded results. Replaces four subsections totalling 195 words.
% Every number verified: NUMBERS.md §1 (main table), §2 (Kendall W), plus
% direct recomputation of the CON baseline count (10/12) and the 60-item
% human subset (Gemini 2.5 Flash 85.0, Gemma 4 E4B 63.3).

\subsection{Overall Performance}
\label{sec:main-results}

Table~\ref{tab:mainresults} reports strict accuracy for all twelve models on the full 406-item
benchmark. Scores span 75.1\% (DeepSeek-R1-Distill) to 89.7\% (Gemini 2.5 Flash), a 14.6-point range
that is wide enough to separate the field but narrow enough that most adjacent pairs are not
individually distinguishable.

The task means are more informative than the model means. Averaged across models, contradiction
detection is highest at 91.0\%, followed by regulatory interpretation at 83.8\%, temporal reasoning
at 77.0\%, and numerical reasoning at 65.5\%. That ordering is stable: Kendall's $W$ across the four
task profiles is 0.853 ($\chi^2 = 30.7$, $p < 0.001$, $n = 12$), so the models largely agree about
which tasks are hard, even where they disagree about their own overall standing. Difficulty here is
a property of the tasks rather than of any particular system.

Two of these numbers should not be read at face value, for opposite reasons. The contradiction
figure is inflated by class imbalance and is addressed in Section~\ref{sec:tasklevel}. The Gemini
2.5 Pro row is depressed by scoring: at 76.1\% strict it sits ninth, below the much smaller Gemini
2.5 Flash, and the gap is almost entirely a formatting artifact of its verbose output style---under
corrected scoring it reaches 95.1\%. It is the same effect documented at length in
Section~\ref{sec:regime}, visible here as a single anomalous row.

\subsection{Task-Level Analysis}
\label{sec:tasklevel}

The four tasks differ sharply in how well they separate systems.

\textbf{Numerical reasoning discriminates most.} Its 35.9-point spread (84.8\% down to 48.9\%) is
the widest of any task, and it is the only task where a frontier model falls to near-chance
performance. Multi-step arithmetic over thresholds embedded in regulatory prose requires correct
extraction \emph{and} correct execution, and failure at either stage produces a wrong answer.

\textbf{Contradiction detection discriminates least, despite the highest mean.} 53 of its 62 items
carry the gold label `No', so a classifier that always answers `No' scores 85.5\%. Against that
baseline the 91.0\% task mean is a 5.5-point margin, not a strong result, and 10 of 12 models
exceed the baseline while Mistral-7B (80.6\%) and Gemma 4 E4B (79.0\%) fall below a trivial
majority classifier. High accuracy on an imbalanced binary task is close to uninformative, and we
report balanced accuracy in Appendix~\ref{app:con} rather than leaning on the raw figure.

\textbf{Regulatory interpretation and temporal reasoning sit in between}, at 20.7-point and
29.5-point spreads. Temporal reasoning is where the top of the table separates: Qwen3-32B leads at
92.3\% while Gemini 2.5 Pro reaches only 64.1\%.

\subsection{Difficulty Stratification}
\label{sec:difficulty}

Items were labelled easy (160), medium (182), or hard (64) at annotation time by the number of
reasoning steps required. If that label tracked model difficulty as intended, accuracy should fall
monotonically from easy to hard. It does, for exactly three of twelve models: Gemini 2.5 Flash
(92.5\% $\to$ 84.4\%), GPT-OSS 120B (79.4\% $\to$ 73.4\%), and most sharply Gemini 2.5 Pro
(83.1\% $\to$ 68.8\%, a 14.4-point decline). The other nine do not decline monotonically at all---
the dominant pattern in this benchmark is not the one the difficulty label was designed to produce.

Two non-monotonic shapes recur. Six models---Qwen3-32B, LLaMA-3.3-70B, Llama 4 Scout 17B, Kimi K2,
LLaMA-3-8B, and DeepSeek-R1-Distill---score \emph{higher} on hard items than on medium ones, and for
LLaMA-3.3-70B this extends further: hard items (90.6\%) outscore even the easy ones (79.4\%), an
11.2-point inversion. Gemma 4 E4B and GPT-OSS 20B show the opposite shape, a dip at medium difficulty
with partial recovery on hard (Gemma: 82.5\% easy $\to$ 74.7\% medium $\to$ 81.2\% hard). Neither
shape is what "harder" was meant to produce.

We read this as evidence about the difficulty label, not about any individual model. Reasoning-step
count was the criterion used to assign difficulty, but hard items in this benchmark are typically
multi-instrument amendment-chain questions that state their cues explicitly in the text, whereas
medium items more often turn on a single subtle qualifying clause---a property closer to
\emph{ambiguity} than to \emph{reasoning depth}. A model with strong long-context tracking can find
an explicit multi-step chain easier to follow than a terse single clause that hinges on one word. We
flag this as a limitation of the annotation scheme (Section~\ref{sec:limitations}) rather than
adjusting the labels post hoc, and note that it strengthens the case for the effective-size analysis
in Section~\ref{sec:effsize}: nominal difficulty and empirical difficulty are visibly different
quantities on this benchmark, in the same way nominal size and discriminative size are.

\subsection{A Human Reference Point}
\label{sec:human}

To anchor the accuracy scale, a single non-specialist evaluator---a university-educated adult with
no professional background in finance or law---answered a 60-item subset drawn from the medium and
hard strata, under the same context-only constraint given to the models and scored by the identical
pipeline.

The evaluator scored 80.0\% (48/60; 95\% Wilson CI [68.2\%, 88.2\%]), with a task profile of
REG 100\% (11/11), TMP 87.5\% (14/16), CON 82.4\% (14/17), and NUM 56.2\% (9/16). On the same 60
items, five of the twelve models score above the evaluator and seven below. The best, Gemini 2.5
Flash, reaches 85.0\% against the evaluator's 80.0\%, a difference that is not statistically
significant (paired bootstrap $p = 0.43$); the weakest, Gemma 4 E4B at 66.7\%, is significantly
worse ($p = 0.020$).

We state the limits of this comparison plainly, because they are severe. It is one evaluator on 60
items. It supports no population-level claim about human performance, and a domain specialist would
almost certainly score higher---particularly on the numerical items, where a finance professional
would recognise the standard calculation forms this evaluator had to derive from the passage. The
comparison is an anchor for the scale, not an estimate of human capability, and we report it only
because the alternative---presenting model accuracies with no human referent at all---leaves the
reader with no way to judge whether 85\% on this benchmark is good.

One asymmetry does survive those caveats and is worth recording. The evaluator was perfect on
regulatory interpretation and worst on numerical reasoning, which is the same ordering the models
produce. The hardest task for this benchmark's models is also the hardest for a careful human
reader, which suggests the difficulty is a property of multi-step quantitative inference over
regulatory prose rather than an artifact of how models are prompted or scored.
